% 封面
\begin{titlepage}
    \thispagestyle{empty}
    
    \definecolor{coverpaper}{RGB}{247, 244, 238}
    \definecolor{covernavy}{RGB}{24, 46, 86}
    \definecolor{covergold}{RGB}{186, 146, 76}
    \definecolor{coverteal}{RGB}{84, 129, 138}
    \definecolor{coverbrick}{RGB}{168, 92, 72}
    \definecolor{coverplum}{RGB}{103, 96, 140}
    \definecolor{coverline}{RGB}{193, 202, 215}

    \begin{tikzpicture}[remember picture, overlay]
        % 背景与结构色块
        \fill[coverpaper] (current page.south west) rectangle (current page.north east);
        \fill[covernavy] (current page.south west) rectangle ([xshift=4.4cm]current page.north west);
        \fill[covergold] ([xshift=4.4cm, yshift=-1.28cm]current page.north west)
            rectangle ([yshift=-1.06cm]current page.north east);

        % 左侧竖向辅助线
        \draw[white!18, line width=0.5pt]
            ([xshift=2.86cm, yshift=2.15cm]current page.south west) --
            ([xshift=2.86cm, yshift=-1.4cm]current page.north west);
        \draw[white!18, line width=0.5pt]
            ([xshift=3.76cm, yshift=2.15cm]current page.south west) --
            ([xshift=3.76cm, yshift=-1.4cm]current page.north west);

        % 左侧信息带
        \node[
            anchor=west,
            rotate=90,
            text=white!74,
            font=\fontsize{11pt}{13pt}\selectfont\bfseries
        ] at ([xshift=1.0cm, yshift=3.0cm]current page.south west)
            {HUNAN NORMAL UNIVERSITY};
        \node[
            anchor=west,
            rotate=90,
            text=covergold!98,
            font=\fontsize{13pt}{16pt}\selectfont\bfseries
        ] at ([xshift=2.35cm, yshift=3.0cm]current page.south west)
            {数学物理方法};
        \node[
            anchor=west,
            rotate=90,
            text=white!58,
            font=\fontsize{8.9pt}{10.8pt}\selectfont
        ] at ([xshift=3.31cm, yshift=3.0cm]current page.south west)
            {复变函数 / 积分变换 / 数理方程 / 特殊函数};

        \node[
            anchor=south west,
            text=white,
            align=left,
            text width=3.05cm,
            font=\fontsize{10.2pt}{12.3pt}\selectfont\bfseries
        ] at ([xshift=0.72cm, yshift=0.98cm]current page.south west)
            {物理与电子\\科学学院};
        \node[
            anchor=south west,
            text=white!68,
            align=left,
            text width=3.05cm,
            font=\fontsize{8.6pt}{10pt}\selectfont
        ] at ([xshift=0.72cm, yshift=0.14cm]current page.south west)
            {School of Physics\\and Electronics};

        % 右上角圆环与坐标感元素
        \begin{scope}[shift={([xshift=-1.55cm, yshift=-2.65cm]current page.north east)}]
            \draw[covergold!55, line width=0.95pt] (0, 0) circle (1.85cm);
            \draw[covergold!18, line width=0.45pt] (0, 0) circle (1.0cm);
            \draw[coverline!72, line width=0.45pt] (-2.3, 0) -- (2.45, 0);
            \draw[coverline!72, line width=0.45pt] (0, -2.15) -- (0, 2.35);
            \node[text=coverline!74!black, font=\fontsize{8.8pt}{10pt}\selectfont] at (2.02, 0.26) {$\operatorname{Re}$};
            \node[text=coverline!74!black, font=\fontsize{8.8pt}{10pt}\selectfont] at (0.36, 2.02) {$\operatorname{Im}$};
        \end{scope}

        % 主标题与书目信息
        \node[
            anchor=west,
            text=covernavy!82,
            font=\fontsize{12.5pt}{15pt}\selectfont
        ] at ([xshift=5.1cm, yshift=-2.95cm]current page.north west)
            {湖南师范大学物理与电子科学学院};
        \node[
            anchor=west,
            text=black!42,
            align=left,
            text width=8.7cm,
            font=\fontsize{10.5pt}{13pt}\selectfont
        ] at ([xshift=5.1cm, yshift=-3.7cm]current page.north west)
            {从物理问题出发，面向本科生自学与课程教学的数学物理方法讲义};

        \node[
            anchor=west,
            text=covernavy,
            font=\fontsize{35pt}{42pt}\selectfont\bfseries
        ] at ([xshift=5.0cm, yshift=-5.55cm]current page.north west)
            {数学物理方法};
        \node[
            anchor=west,
            text=covergold!92!black,
            font=\fontsize{12.5pt}{16pt}\selectfont\itshape
        ] at ([xshift=5.1cm, yshift=-6.85cm]current page.north west)
            {Mathematical Methods for Physics};

        \draw[covergold, line width=0.9pt]
            ([xshift=5.1cm, yshift=-7.65cm]current page.north west) --
            ([xshift=12.2cm, yshift=-7.65cm]current page.north west);

        \node[
            anchor=west,
            rounded corners=2pt,
            draw=covergold,
            fill=covergold!10,
            inner xsep=10pt,
            inner ysep=6pt,
            text=covergold!85!black,
            font=\fontsize{12pt}{14pt}\selectfont\bfseries
        ] at ([xshift=5.1cm, yshift=-8.72cm]current page.north west)
            {课程讲义};

        \begin{scope}[shift={([xshift=5.1cm, yshift=-10.18cm]current page.north west)}]
            \foreach \x/\col/\txt in {
                0.0/coverbrick/复变函数,
                2.8/coverteal/积分变换,
                5.95/covergold/数理方程,
                8.95/coverplum/特殊函数
            }{
                \node[
                    anchor=west,
                    rounded corners=2pt,
                    draw=\col!65,
                    fill=\col!9,
                    inner xsep=8pt,
                    inner ysep=4pt,
                    text=covernavy!88,
                    font=\fontsize{10.1pt}{12pt}\selectfont
                ] at (\x, 0.18) {\txt};
            }
        \end{scope}

        \node[
            anchor=west,
            text=black!56,
            font=\fontsize{10.5pt}{13pt}\selectfont
        ] at ([xshift=5.1cm, yshift=-11.42cm]current page.north west)
            {物理学、光电信息科学等专业本科生适用};

        \node[
            anchor=west,
            text=covernavy,
            font=\fontsize{17pt}{20pt}\selectfont\bfseries
        ] at ([xshift=5.1cm, yshift=-12.82cm]current page.north west)
            {陈祖成};
        \node[
            anchor=west,
            text=black!50,
            font=\fontsize{10.5pt}{12pt}\selectfont
        ] at ([xshift=7.72cm, yshift=-12.84cm]current page.north west)
            {编著};

        % 下半部分主视觉
        \begin{scope}[shift={([xshift=-5.15cm, yshift=8.55cm]current page.south east)}]
            \draw[coverline, line width=0.55pt] (-4.9, 0) -- (4.9, 0);
            \draw[coverline, line width=0.55pt] (0, -4.3) -- (0, 4.3);

            \draw[coverplum!38, line width=0.55pt, dash pattern=on 2.5pt off 2.2pt]
                (0, 0) circle (4.15);
            \draw[covergold!80, line width=1.15pt] (0, 0) circle (3.15);
            \draw[covergold!36, line width=0.55pt] (0, 0) circle (1.55);

            \fill[covergold!14] (0, 0) -- (1.05, 0)
                arc[start angle=0, end angle=57, radius=1.05] -- cycle;
            \draw[covergold!90!black, line width=0.9pt]
                (1.05, 0) arc[start angle=0, end angle=57, radius=1.05];

            \draw[
                covernavy,
                line width=1.2pt,
                -{Stealth[length=2.9mm, width=1.9mm]}
            ] (0, 0) -- (57:3.15);
            \node[
                text=covernavy,
                font=\fontsize{10.5pt}{12pt}\selectfont
            ] at (57:3.7) {$re^{\mathrm{i}\theta}$};
            \node[
                text=covergold!85!black,
                font=\small
            ] at (1.0, 0.55) {$\theta$};

            \foreach \y/\col/\lw in {
                -2.65/coverbrick!78/0.95,
                -1.35/coverteal!92/1.0,
                0/covernavy!32/0.8,
                1.35/coverteal!92/1.0,
                2.65/coverbrick!78/0.95
            }{
                \draw[\col, line width=\lw pt, samples=240, domain=-4.9:4.9]
                    plot (\x, {\y + 0.17*sin(58*\x) + 0.1*sin(150*\x)});
            }

            \node[
                text=black!42,
                align=left,
                font=\fontsize{10pt}{12pt}\selectfont
            ] at (-2.35, -3.55)
                {$\displaystyle \frac{\partial^2 u}{\partial t^2}=c^2\nabla^2u$};
        \end{scope}
    \end{tikzpicture}
    
\end{titlepage}
